\section{Introduction}%\label{sec1}

The class of nonassociative algebras with one binary operation satisfying the identities of left symmetry
%
\begin{equation}\label{eq:LSym}
(xy)z - x(yz) = (yx)z - y(xz)
\end{equation}
%
and right commutativity
%
\begin{equation}\label{eq:RCom}
(xy)z = (xz)y
\end{equation}
%
is known as the variety of {\em Novikov algebras}.
Relations \eqref{eq:LSym} and \eqref{eq:RCom} emerged
in \cite{BN1985, GD79} as a~tool
for expressing certain conditions on the components of
a tensor of rank~3 considered
as a collection of structure constants of a finite-dimensional
 algebra with one bilinear operation.
 In~\cite{GD79}, it was a sufficient condition
for a differential operator to be Hamiltonian;
in \cite{BN1985} it was a condition to guarantee
the Jacobi identity for a generalized Poisson bracket in
the framework of Hamiltonian formalism for partial differential
equations of hydrodynamic type.

Novikov algebras may be obtained from (associative) commutative algebras
with a derivation by means of the following operation-transforming functor (see \cite{GD79}).
Assume $A$ is a commutative algebra
with a multiplication~$*$, and let $d$ be a derivation of $A$, i.e.,
 a linear operator $d\colon A\to A$
satisfying the Leibniz rule
\[
d(a*b) = d(a)*b+a*d(b), \qquad a,b\in A.
\]
Then the new operation
\begin{equation}\label{eq:Gelfand}
ab = a*d(b), \qquad a,b\in A,
\end{equation}
satisfies the identities \eqref{eq:LSym} and \eqref{eq:RCom}.

In \cite{DzhLofwall2002}, it was proved that the free Novikov
algebra $\Nov \langle X\rangle $ generated by a set $X$ is
a subalgebra of the free differential commutative algebra
with respect to the operation~\eqref{eq:Gelfand}.
Moreover, it was shown in \cite{BCZ-2017} that every Novikov algebra
can be embedded into an appropriate
commutative algebra with a derivation.

One may generalize the relation between commutative algebras with a derivation and Novikov algebras as follows.
Let $\Var$ be a class of all algebras over a field $\Bbbk $ with one or more bilinear operations which is closed under direct products, subalgebras, and homomorphic images (HSP-class or {\em variety}).
By the classical Birkhoff's theorem, a variety consists of all algebras
that satisfy some family of identities.
We will assume that $\Var $ is defined by multi-linear identities (this is not a restriction if $\operatorname{char}\Bbbk =0$).
For every $A\in \Var $, let $\End(A)$ stand for the set
of all linear operators on the space~$A$.
Then a linear operator $d\in \End(A)$
is called a derivation if the analogue of the Leibniz rule
holds for every binary operation in~$A$.
The set of all derivations of a given algebra $A$ forms a subspace (even a Lie subalgebra) $\Der(A)$ of $\End(A)$.
The class of all pairs $(A,d)$, $A\in \Var$, $d\in \Der(A)$,
is also a variety denoted $\Var\Der$ defined by multi-linear identities.
Given $(A,d)\in \Var\Der$, denote by
$A^{(d)}$ the same space $A$ equipped with two bilinear operations
$\prec$, $\succ $ for each operation on~$A$:
\[
a\prec b = ad(b),\qquad a\succ b = d(a)b, \qquad a,b\in A.
\]
The class of all systems $A^{(d)}$, $(A,d)\in \Var\Der$,
is closed under direct products, so all homomorphic images of
all their subalgebras form a variety denoted $D\Var$,
called the {\em derived variety} of $\Var $ (see \cite{KSO2019}).
Alternatively, $D\Var $ consists of all algebras with duplicated family
of operations that satisfy all those identities that hold on all $A^{(d)}$
for all $(A,d)\in \Var\Der$.

For example, if $\Com $ is the variety of commutative (and associative) algebras then $D\Com = \Nov$ since
$x\prec y = y\succ x$. In general, the description of
$D\Var $ may be obtained in the language of operads and their
Manin products \cite{KSO2019}. Namely, if we identify the notations
for a variety and its governing operad \cite{GK-1994}
then the operad
$D\Var $ coincides with the Manin white product of operads
$\Var $ and $\Nov $.

As a corollary, the free algebra $D\Var \langle X\rangle $
in the variety $D\Var $ generated by a set $X$
is isomorphic to
the subalgebra in the free algebra in the variety
$\Var\Der$ generated by~$X$. However, it is not clear that
the following {\em embedding statement} holds in general:
 every
$D\Var$-algebra can be embedded into an appropriate
differential $\Var$-algebra (or $\Var\Der$-algebra).
In other words, the problem is to determine whether the
class of all subalgebras of all algebras $A^{(d)}$,
$(A,d)\in \Var\Der$,
is closed under homomorphic images.
Positive answers were obtained
for $\Var =\Com $ \cite{BCZ-2017}, $\Var=\Lie$ \cite{KSO2019}, $\Var=\Perm$ \cite{KS2022}, $\Var=\As$ \cite{SK-2022}.

In this paper, we derive a sufficient condition for a positive
answer to the embedding statement for a given $\Var $.
Namely, if the Manin white product $\Nov\circ \Var$ of the corresponding operads coincides with the Hadamard product
$\Nov\otimes \Var $ then the embedding statement holds for $\Var$.

We also find an example of a variety $\Var $ governed by
a binary quadratic operad such that the embedding
statement fails for $\Var$. It turns out that
the variety of Zinbiel algebras $\Zinb $ (also known as commutative dendriform algebras, pre-commutative algebras, dual Leibniz algebras, half-shuffle algebras) introduced in \cite{Loday} provides such an example:
there exists a $D\Zinb $-algebra which cannot be embedded into a Zinbiel algebra with a derivation.
To our knowledge, this is the first example of such a variety $\Var$ that an algebra from $D\Var$ does not embed into a~differential $\Var$-algebra.

The operad $\Zinb $ governing the variety of Zinbiel algebras
is a particular case of a general construction called {\em dendriform splitting} of an operad \cite{BBGN}.
For every binary operad $\Var $ (not necessarily quadratic,
see, e.g., \cite{GubKol2014}) there exists an operad
$\pre\Var $ governing the class of systems with duplicated
family of operations. The generic example of a $\pre\Var$ algebra may be obtained from a $\Var$-algebra
with a Rota--Baxter operator $R$ of weight zero
(see, e.g., \cite{LiGuo_RBABook}).
If $A\in \Var$ and $R\colon A\to A$ is such an operator then
$(A,\vdash, \dashv)$, where
\[
a\vdash b = R(a)b,\qquad a\dashv b = aR(b), \qquad a,b\in A,
\]
is a $\pre\Var $-algebra.
In this context,
$\pre\Com =\Zinb$ (since $a\vdash b = b\dashv a$),
$\pre\Lie$ is the classical variety of left-symmetric algebras (relative to $\vdash $), $\pre\As$ is exactly the variety of dendriform algebras \cite{Loday}.

The theory of pre-algebras and relations between them is similar
in many aspects to the theory of ``ordinary'' algebras.
For example, every $\pre\As $-algebra with respect to the ``dendriform commutator'' $a\vdash b - b\dashv a$ is a $\pre\Lie$-algebra. The properties of the corresponding left adjoint functor (universal envelope) are close to what we have
for ordinary Lie algebras
\cite{DotsTamaroff, Gubarev}.
The class of pre-Novikov algebras has been recently studied in \cite{Hong2023}: it coincides with $D\Zinb $.
Therefore, our results show that the embedding statement
cannot be transferred from ordinary algebras to pre-algebras.


\section{Derived algebras and Manin white product of binary operads}

The details about (symmetric) operads may be found, for example,
in \cite{Bremner-Dotsenko}.
For an operad $\mathcal P$ denote
$\mathcal P(n)$ the linear space (over a base field $\Bbbk $)
of degree $n$ elements of $\mathcal P$, the action of a~permutation $\sigma \in S_n$
on an element $f\in \mathcal P(n)$ is denoted $f^\sigma $,
let $\id \in \mathcal P(1) $ stand for the identity element,
and the composition rule
\[
\mathcal P(n)\otimes \mathcal P(m_1)\otimes \dots \otimes \mathcal P(m_n) \to \mathcal P(m_1+\dots + m_n)
\]
is denoted $\gamma^{m_1+\dots+m_n}_{m_1,\dots , m_n}$.

Recall that an operad $\mathcal P$ is said to be binary, if
$\mathcal P(1)=\Bbbk \mathrm{id}$ and the entire
$\mathcal P$ is generated (as an operad) by its degree 2 space
$\mathcal P(2)$.

Let us fix a binary operad $\mathcal P$.
For every
linear space $A$, consider an operad
$\End_A$
with $\End_A(n)=\mathrm{End}\,(A^{\otimes n}, A)$.
A~morphism of operads
$\mathcal P \to \End_A$ defines an algebra structure
on $A$ with a set of binary operations corresponding
to the generators of $\mathcal P$ from $\mathcal P(2)$.
The class of all such algebras is a variety
 of $\mathcal P$-algebras
defined by multi-linear identities corresponding
to the defining relations of the operad~$\mathcal P$.

Conversely, every variety of algebras with binary operations
defined by multi-linear identities gives rise to
an operad $\mathcal P$ constructed in such a way that
$\mathcal P(n)$ is the space of multi-linear elements
of degree~$n$ in the variables $x_1,\ldots, x_n$
of the free algebra in this variety generated by the countable set
$\{x_1,x_2,\ldots \}$
(see, e.g., \cite[Section~1.3.5]{Kol2008_smz} for details).
In particular, $\id\in \mathcal P(1)$
is presented by the element $x_1$.
Then the variety under consideration consists
exactly of all $\mathcal P$-algebras, in other words, it is
{\em governed} by the operad~$\mathcal P$.

\begin{Example}
Let $\mu,\nu \in \mathcal P(2)$. Suppose we have identified
$\mu $ with $x_1x_2$ and $\nu $ with $x_1*x_2$.
Then{\samepage
\begin{gather*}
\begin{split}
&\gamma^3_{1,2}(\mu, \id, \nu )
= \gamma^3_{1,2} (x_1x_2, x_1, x_1*x_2)
= x_1(x_2*x_3),
\\
&\gamma^4_{2,2}\bigl(\nu^{(12)}, \mu, \nu \bigr)
= \gamma^4_{2,2}(x_2*x_1, x_1x_2, x_1*x_2)
= (x_3*x_4)*(x_1x_2),
\end{split}
\end{gather*}
and so on.}
\end{Example}

We will not distinguish notations for an operad $\mathcal P$
and for the corresponding variety of $\mathcal P$-algebras.

\begin{Example}
Suppose $\mathcal F_2$ is the free binary operad with
$\mathcal F_2(2) \simeq \Bbbk S_2$ (as symmetric modules).
Then the class of $\mathcal F_2$-algebras
coincides with the variety of all nonassociative
algebras with one non-symmetric binary operation,
i.e., $\mathcal F_2(n)$ may be identified with the
linear span of all bracketed monomials
$\bigl(x_{\sigma(1)}\cdots x_{\sigma(n)}\bigr)$, $\sigma \in S_n$,
so $\dim \mathcal F_2(n) = n! C_{n-1}$, where $C_k$ is the $k$th
Catalan number.
\end{Example}

If $\mathcal P$ is an operad governing a variety
of algebras with one binary operation then
$\mathcal P$ is a~homomorphic image of $\mathcal F_2$.
If the kernel of the morphism is generated (as an operadic ideal) by elements from $\mathcal F_2(3)$ then
the operad $\mathcal P$ is said to be {\em quadratic}.
The same definition works for operads with more than one generator.

\begin{Example}\label{exmp:Zinbiel}
Let $\Zinb$ stand for the variety of algebras with one multiplication
satisfying the identity
%
\begin{equation}\label{eq:Zinbiel}
(x_1\cdot x_2)\cdot x_3
= x_1\cdot (x_2\cdot x_3 + x_3\cdot x_2),
\end{equation}
%
known as the Zinbiel identity \cite{Loday}.
Then the space $\Zinb(n)$, $n\ge 1$, is spanned by
linearly independent monomials
\[
\bigl( x_{\sigma^{-1}(1)}\cdot \bigl(x_{\sigma^{-1}(2)}\cdot \bigl(\cdots \bigl(x_{\sigma^{-1}(n-1)}\cdot x_{\sigma^{-1}(n)}\bigr)\cdots \bigr)\bigr) \bigr)
=
[x_{1}x_{2}\cdots x_{n-1}x_{n}]^\sigma ,
\qquad \sigma \in S_n,
\]
so $\dim \Zinb(n)=n!$.

The product of two right-normed words in a Zinbiel algebra
can be expressed via shuffle permutations $S_{n,m}\subset S_{n+m}$: if $u= x_1\cdots x_n$, $v = x_{n+1}\cdots x_{n+m}$
then
\[
[u]\cdot [v]
=\sum\limits_{\sigma\in S_{n,m}\colon \sigma (1)=1}
[uv]^\sigma .
\]
The new product $[u]\cdot [v] + [v]\cdot [u]$
on a Zinbiel algebra
is associative and commutative (it is known as the {\em shuffle product}
of words).
\end{Example}


Suppose $\mathcal P_1$ and $\mathcal P_2$ are two binary operads.
Then the {\em Hadamard product} of $\mathcal P_1$ and $\mathcal P_2$ is the operad denoted
$\mathcal P= \mathcal P_1\otimes \mathcal P_2$ such that
$\mathcal P(n)= \mathcal P_1(n)\otimes \mathcal P_2(n)$, $n\ge 1$,
the action of $S_n$ on $\mathcal P(n)$ and the composition rule are defined in
the componentwize way.

The operad $\mathcal P_1\otimes \mathcal P_2$ may not be a binary one (in this paper, we will deal
with such an example below). The sub-operad of $\mathcal P_1\otimes \mathcal P_2$ generated
by $\mathcal P_1(2)\otimes \mathcal P_2(2)$ is
called the {\em Manin white product} of $\mathcal P_1$ and $\mathcal P_2$ denoted by
$\mathcal P_1\circ\mathcal P_2$.
(For some operads $\mathcal P_1$, it may happen that
$\mathcal P_1\circ \mathcal P_2 = \mathcal P_1\otimes \mathcal P_2$
for all $\mathcal P_2$; in \cite{Vallette2008}, all such operads were described.)

If $\mathcal P_1$ and $\mathcal P_2$ are quadratic binary operads then so is
$\mathcal P_1\circ \mathcal P_2$. The defining relations of the last operad can be found
as follows \cite{GK-1994}. Suppose $R_i\subseteq \mathcal P_i(3)$,
$i=1,2$, are the spaces of defining relations of
$\mathcal P_i$ presented in the form of multi-linear identities in $x_1$, $x_2$, $x_3$.
Consider the space~$E(3)$ spanned by all possible compositions of degree~3 of operations from
$\mathcal P_1(2)\otimes \mathcal P_2(2)$.
Then the defining relations of $\mathcal P_1\circ \mathcal P_2$
form the space $E(3)\cap (\mathcal P_1(3)\otimes R_2 + R_1\otimes \mathcal P_2(3))$.

\begin{Example}\label{exmp:Nov-Zinb}
Let $\mathcal P_1=\Nov$, $\mathcal P_2=\Zinb$.
Then $\mathcal P_1(2)\otimes \mathcal P_2(2)$ is spanned by four elements
\begin{alignat*}{3}
& x_1\prec x_2 = x_1x_2\otimes x_1x_2,
 \qquad &&
 x_1\succ x_2 = x_2x_1\otimes x_1x_2, &
 \\
& x_2\prec x_1 = x_2x_1\otimes x_2x_1,
 \qquad &&
 x_2\succ x_1 = x_1x_2\otimes x_2x_1. &
\end{alignat*}
In order to find $E(3)$, calculate all monomials of degree 3 in $x_1$, $x_2$, $x_3$
with operations $\succ$, $\prec$. For example,
%
\begin{align*}
 (x_1\succ x_3)\prec x_2 &= \gamma^3_{1,2} (x_2\prec x_1, \id, x_1\succ x_2)^{(12)}
 \\
 &=
 \gamma^3_{1,2}(x_2x_1, \id, x_2x_1)^{(12)}\otimes \gamma^3_{1,2}(x_2x_1, \id, x_1x_2)^{(12)}
 \\
& =
(x_3x_2)x_1^{(12)}\otimes (x_2x_3)x_1^{(12)}
=
(x_3x_1)x_2\otimes (x_2x_3)x_1 \in \Nov(3)\otimes \Zinb(3).
\end{align*}
%
In the same way, the expressions for all 48 monomials may be calculated
in $\Nov(3)\otimes \Zinb(3)$. In order to get defining relations
of $\Nov\circ \Zinb$ it is enough to find those linear combinations
of these monomials that are zero in $\Nov(3)\otimes \Zinb(3)$.
This is a routine problem of linear algebra.
As a result, we obtain the following identities:
%
\begin{gather}
(x_1\prec x_2)\prec x_3=(x_1\prec x_3)\prec x_2,\nonumber\\
x_1\succ (x_2\succ x_3) = (x_1\succ x_3)\prec x_2-x_1\succ (x_3\prec x_2),\nonumber\\
(x_1\succ x_2)\succ x_3=(x_1\succ x_3)\succ x_2 +(x_1\succ x_3)\prec x_2-(x_1\succ x_2)\prec x_3, \label{eq:pre-Novikov}\\
%
(x_1\prec x_2)\succ x_3=x_1\prec (x_2\succ x_3)+x_1\prec (x_3\prec x_2)
 +(x_1\succ x_3)\prec x_2-(x_1\prec x_3)\prec x_2.\nonumber
\end{gather}
\end{Example}

\begin{Remark}%\label{rem:Xu-GD}
Novikov algebras are closely
related to conformal algebras in the sense
of \cite{Kac:VAforbeginners}.
Namely, if $V$ is a Novikov algebra then
the free $\Bbbk [\partial ]$-module $C=\Bbbk[\partial ]\otimes V$
equipped with a~sesqui-linear $\lambda $-product
\[
[u{}_{(\lambda )} v] = \partial (vu) + \lambda (vu + vu),
\qquad u,v\in V,
\]
is a Lie conformal algebra \cite{Xu1999}. The reason of this
relation is that the expression for the generalised Poisson
bracket in \cite{BN1985} has the same form as the expression
for the commutator of local chiral fields in \cite{Kac:VAforbeginners}.
\end{Remark}

\begin{Remark}\label{rem:Conformal}
The variety governed by the operad $\Nov\circ \Zinb $
is also related to conformal algebras.
It is straightforward to check that if $V$ is an algebra over a field $\Bbbk $, $\operatorname{char}\Bbbk =0$, with two
operations~$\prec$,~$\succ $ satisfying the identities
\eqref{eq:pre-Novikov}
then the free
$\Bbbk [\partial ]$-module $C=\Bbbk[\partial ]\otimes V$
equipped with a~sesqui-linear $\lambda $-product
\[
\bigl(u{}_{(\lambda )} v\bigr) = \partial (v\prec u) + \lambda (v\succ u + v\prec u),
\qquad u,v\in V,
\]
is a left-symmetric conformal algebra \cite{HongLi2015}.
We will explain this relation in the last section.
\end{Remark}

Let $\Var $ be a binary operad, we use the same notation for the
corresponding variety of algebras.
Given an algebra $A$ in $\Var $, denote by $\Der(A) $
the set of all derivations of $A$. Recall that
a derivation of $A$ is a linear map $d\colon A\to A$
such that
\[
 d(\mu(a,b)) = \mu(d(a),b) + \mu(a,d(b)), \qquad a,b\in A,
\]
for all operations $\mu $ from $\Var (2)$.
For a derivation $d$ of an algebra $A$, denote by
$A^{(d)}$ the linear space $A$ equipped with {\em derived}
operations
%
\begin{equation}\label{eq:Der-op}
 \mu^{\succ }(a,b) = \mu(d(a),b), \quad \mu^{\prec }(a,b) = \mu(a,d(b)), \qquad a,b\in A,
\end{equation}
%
for all $\mu $ in $\Var (2)$.
The variety generated by the class of systems $A^{(d)}$ for all $A\in \Var$ and $d\in \Der(A)$
is denoted by $D\Var$, the {\em derived variety} of $\Var $.
Obviously, the class of all such $A^{(d)}$ is closed under direct products,
so, in order to get $D\Var$, one has to consider all homomorphic images of all subalgebras of these $A^{(d)}$.
In many cases, it is enough to consider just subalgebras.
The problem is to decide whether homomorphic images are actually needed.

\begin{Theorem}[\cite{KSO2019}]
For a binary operad $\Var $, the variety $D\Var $ is governed
by the operad ${\Nov\circ \Var}$.
\end{Theorem}

For example,
all relations that hold
on every Zinbiel algebra with a derivation relative to
the operations $a\succ b = d(a)b$, $a\prec b = ad(b)$
follow from the identities~\eqref{eq:pre-Novikov}.

If $\Var = \Com$ is the variety of associative and commutative algebras
then $D\Var = \Nov$, as follows from the construction of the free Novikov
algebra~\cite{DzhLofwall2002}. Here we have to mention
that commutativity implies $a\succ b = b\prec a$ in every
algebra from $D\Com $.

If $\Var = \Lie $ then the algebras from $D\Var $ form exactly the class of
all $\mathcal F_2$-algebras with one binary operation $a\succ b = -b\prec a$ \cite{KSO2019}.
Note that $\dim \Nov(n) = \binom{2n-2}{n-1}$ \cite{Dzhumad-codimgrowth}, and $\Lie(n)=(n-1)!$.
Hence,
$\dim (\Nov\otimes \Lie)(n) = \frac{(2n-2)!}{(n-1)!}$ which is equal
to $n!C_{n-1}$, where $C_n$ is the $n$th Catalan number.
The number $n!C_{n-1}$ coincides with the $n$th dimension
of the operad~$\mathcal F_2$.
Hence,
$\Nov\circ \Lie = \Nov \otimes \Lie $.

Suppose $\Var $ is a binary operad, $D\Var = \Nov\circ \Var$,
$D\Var\langle X\rangle $ is the free $D\Var$-algebra generated by a countable set
$X=\{x_1,x_2,\dots \}$. Denote by $F=\Var\Der\langle X,d\rangle $ the
free differential $\Var $-algebra generated by $X$ with one derivation $d$.
Then there exists a homomorphism $\tau \colon D\Var\langle X\rangle \to F^{(d)}$
sending $X$ to $X$ identically.
An element from $\ker \tau $ is an identity that holds on all $\Var $-algebras
with a derivation relative to the derived operations \eqref{eq:Der-op}.
Hence, $\tau $ is injective, i.e., the free $D\Var$-algebra
can be embedded into the free differential $\Var $-algebra.

The next question is whether every $D\Var$-algebra can be embedded into an
appropriate differential $\Var $-algebra.
This is the same as to decide if the class of all subalgebras of $A^{(d)}$,
$(A,d)\in \Var\Der$, is closed under homomorphisms.
The answer is positive
for $\Var = \Com $ \cite{BCZ-2017}, $\Lie $ \cite{KSO2019}, $\Perm $ \cite{KS2022}, and $\As $ \cite{SK-2022}.
In the following sections we derive a sufficient condition
for $\Var $ to guarantee a positive answer. This condition is
not necessary, but we find an example when the answer is negative.


\section{The weight criterion and special derived algebras}

Let $\Var $ be a binary operad.
An algebra $V$ with two binary operations $\prec$, $\succ $ from the variety
$D\Var $ is called {\em special} if it can be embedded into
a $\Var $-algebra $A$ with a derivation $d$ such that
$u\prec v = ud(v)$ and $u\succ v = d(u)v$ in $A$ for all
$u,v\in V$.
The class of all $\Var $-algebras with a derivation is a variety
since it is defined by identities. The free differential $\Var $-algebra
$\Var\Der\langle X,d\rangle $
generated by a set $X$ is isomorphic as a $\Var$-algebra
to the free $\Var$-algebra $\Var\bigl\langle X^{(\omega )}\bigr\rangle $
generated by the set
\[
X^{(\omega )} = \bigl\{x^{(n)} \mid x\in X,\, n\in \mathbb Z_+ \bigr\},
\]
with the derivation $d$ defined by $d\bigl(x^{(n)}\bigr) = x^{(n+1)}$,
$x\in X$, $n\in \mathbb Z_+$.

For a nonassociative monomial $u\in X^{(\omega )}$ define its
weight $\wt(u)\in \mathbb Z$ as follows. For a single letter
$u=x^{(n)}$, set $\wt(u)=n-1$. If $u=u_1u_2$ then
$\wt(u) = \wt(u_1)+\wt(u_2)$.
Since the defining relations of $\Var\bigl\langle X^{(\omega )}\bigr\rangle $
are weight-homogeneous, we may define the weight function
on $\Var\Der\langle X,d\rangle $.
Note that if $f\in \Var\bigl\langle X^{(\omega )}\bigr\rangle $
is a weight-homogeneous polynomial then
$\wt d(f) = \wt(f)+1$.

\begin{Lemma}\label{lem:Weight-criterion}
 Let $\Var $ be a binary operad such that
 $\Nov\circ \Var = \Nov \otimes \Var $.
Then for every set~$X$ an element $f\in \Var\bigl\langle X^{(\omega )}\bigr\rangle $
belongs to $D\Var \langle X\rangle $ if and only if
$\wt (f)=-1$.
\end{Lemma}

\begin{proof}
The ``only if'' part of the statement does not depend on the particular operad $\Var $.
Indeed,
every formal expression in the variables $X$ relative to binary operations $\prec$
and $\succ $ turns into a weight-homogeneous polynomial of weight $-1$ in $\Var\bigl\langle X^{(\omega )}\bigr\rangle $.

For the ``if'' part, assume $u$ is a monomial of weight $-1$ in the variables $X^{(\omega )}$.
In the generic form,
\[
 u = \bigl(x_{i_1}^{(s_1)}\cdots x_{i_n}^{(s_n)}\bigr), \qquad x_{i_j}\in X,\quad s_j\ge 0,
\]
with some bracketing. Here $s_1+\dots +s_n = n-1$.
Consider the element
\[
[u] = x_{1}^{(s_1)}\cdots x_{n}^{(s_n)} \otimes (x_1\cdots x_n) \in \Nov(n)\otimes \Var(n).
\]
Here the first tensor factor is a differential commutative monomial of degree $n$ and of weight~$-1$
which belongs to $\Nov(n)$ by \cite{DzhLofwall2002}. In the second factor, we put the
nonassociative multi-linear word
obtained from $u$ by removing all derivatives and consecutive
re-numeration of variables (the bracketing remains the same as in~$u$).

By the assumption, $[u]$ belongs to $(\Nov\circ \Var)(n)$, i.e., can be obtained
from $x_1\prec x_2$ and $x_1\succ x_2$ by compositions and symmetric groups actions.
Hence,
the monomial $\bigl(x_{1}^{(s_1)}\dots x_{n}^{(s_n)}\bigr)$ may be expressed in terms of
operations $\succ $ and $\prec$ on the variables $x_1,\dots,x_n\in X$
in the differential algebra $\Var\bigl\langle X^{(\omega )}\bigr\rangle $.
It remains to make the substitution
$x_{j}\to x_{i_j}$ to get the desired expression for $u$ in $D\Var \langle X\rangle $.
\end{proof}
